\documentclass[a4paper, 11pt]{article} % das Papierformat zuerst
\usepackage[utf8]{inputenc}
\usepackage{geometry}
\geometry{a4paper, top=25mm, left=35mm, right=25mm, bottom=25mm}
\usepackage{graphicx}
\usepackage{color}
\usepackage{epstopdf}
\usepackage[T1]{fontenc}
\usepackage{setspace}
\usepackage{textcomp}
\usepackage{tabularx}
\usepackage{blindtext}
\usepackage[ngerman]{babel} %deutsche Silbentrennung
\usepackage{titlesec}
\usepackage{enumitem} 
\usepackage{ifthen}
\usepackage[colorlinks=true,linkcolor=fsi]{hyperref}

\definecolor{fsi}{RGB}{0,81,150}

\newcounter{summe}
\newcommand{\abstimmung}[4]{
	\setcounter{summe}{#3}
	\addtocounter{summe}{#4}
	\begin{flushleft}
		#1\\
	Es wird über den Antrag abgestimmt.
	\end{flushleft}
	\ifthenelse{\equal{#3}{0}\AND\equal{#4}{0}}{
	\begin{center}
		\textbf{Der Antrag wird einstimmig angenommen.}
	\end{center}
	}{
	\begin{center}
		#2 \ifthenelse{\equal{#2}{1}}{Stimme}{Stimmen} dafür, #3 \ifthenelse{\equal{#3}{1}}{Stimme}{Stimmen} dagegen, #4 \ifthenelse{\equal{#42}{1}}{Enthaltung}{Enthaltungen}\\
		\ifthenelse{#2>\value{summe}}{
		\textbf{Der Antrag ist somit angenommen.}
		}{
		\textbf{Der Antrag ist somit abgelehnt.}
		}
	\end{center}
	}  
}


% VARIABLEN
\newcommand{\protokoller}{Philipp Dippel}
\newcommand{\dateOfMeeting}{29.07.2015}
\newcommand{\TeXer}{Philipp Dippel}
\newcommand{\fsiPresident}{Georg Schäfer}
\newcommand{\protocolID}{9}

\begin{document}
%deckblatt start

\doublespacing
\thispagestyle{empty}

\begin{center}
\includegraphics[width=10.0cm]{../logo_fsi.eps}

\vspace*{\fill}
{\LARGE \textbf{Protokoll der Sitzung des Fachschaftsrates \\vom \dateOfMeeting}}\\
Fachschaftsrat Informatik\\
Trier University of Applied Sciences\\
\vspace{2.5cm}
\textit{
	Protokollführer: \textbf{\protokoller} \\
	\LaTeX - Umsetzung von \TeXer\\
	am \today\\
}
Protokoll ID: \textbf{\#\protocolID}
\vfill
\end{center}

\hspace*{-35cm}
\textcolor{fsi}{\rule{64.9cm}{15pt}}
\pagebreak
 
\setcounter{tocdepth}{2}
\tableofcontents 
\pagebreak

\section{Eröffnung}
Als Protokollführer wird \protokoller~bestimmt.\\
Der Sprecher des Fachschaftsrates \fsiPresident~eröffnet die Sitzung um 19:17 Uhr.
\\\\
\textbf{Es wird festgestellt, dass Tobias Meier, Fabio Gimmillaro und Michael Ochmann nicht zur Sitzung erschienen sind. Marco Lochen tritt der Sitzung um 19:21 bei.}
\\\\
\textbf{Der Fachschaftsrat ist somit beschlussfähig.}
\\\\
Das Protokoll der letzten Sitzung wurde verlesen und genehmigt.

%\abstimmung{Es wird der Antrag gestellt, Dinge zu tun}{5}{6}{2} -- {dafür}{dagegen}{enthaltungen}

\pagebreak

\section{Events}
Die Eventmanager verkünden, dass es keine nennenswerten Neuigkeiten von vergangenen Events gibt und keine weiteren Events in näherer Zeit anstehen. 

\section{Finanzen}
Der Kassenreferent Tim Grundmanns und sein Stellvertreter Philipp Dippel verkünden die Endgeldbestände der Netzwerkinstallationsveranstaltung vom 26. Juli 2015. \\ Es ist mit bedauern zu Verkünden, dass nach Abzug der Getränkekosten vom Lieferanten Herrestahl, die Fachschaft Informatik ein Malus von ca. 100,00 Euro vorzuweisen hat. Kalkulierbar wäre nach Stand der Verkaufszahlen ein Kassenstand von ungefähr +100,00 Euro und die negative Bilanz der Kasse ließe sich nach Aussage der Kassenreferenten nur daher erklären, dass es bereits an der NIV zu einem unerlaubtem zugriff in der Kasse kam. Woher dieser Zugriff kam, lässt sich im Nachhinein nicht mehr Nachweisen. Ein Zugriff auf die  Kasse nach der Veranstaltung kann aufgrund der Tatsache das die Kasse für die Öffentlichkeit unzugänglich gemacht und bewacht wurde ausgeschlossen werden.\\\\ Die Fachschaft ist sich dem Umständen entsprechend daher einig, dass auf Zukünftigen Veranstaltung mit Verkauf eine verschärfte Sicherheit beim Umgang mit dem Kassengeschäft von Nöten ist. 
\\\\
Weitergehend wird erwähnt das unser Kassenreferent Tim Grundmanns in den nächsten StuPa gewählt wurde und somit laut Satzung der Fachschaft für Informatik, bis zur konsolidierenden Sitzung des StuPa, von seinem Amt zurücktreten und aus dem Fachschaftsrat austreten muss.  
\\\\ \textbf{Aufgrund dessen tritt Tim Grundmanns von seinem Amt als Kassenreferent zurück. Ebenso tritt auch Philipp Dippel von seinem Amt als Stellvertreter zurück.}
\\

\pagebreak
\noindent Darauffolgend wird eine öffentliche Wahl eines neuen Kassenreferent und Stellvertreters abgehalten. 

\abstimmung{Zur Wahl als Kassenreferent stellt sich Philipp Dippel zur Verfügung.}{10}{0}{0} %-- {dafür}{dagegen}{enthaltungen}

\abstimmung{Zur Wahl als Stellvertreter stellt sich Elias Broschin zur Verfügung.}{10}{0}{0}% -- {dafür}{dagegen}{enthaltungen}

Beide nehmen ihr Amt an. 

\section{Raum}
Der Raumwart Benjamnin Albsmeier verkündet, dass die Putzzeiten mittlerweile besser als zuvor eingehalten werden, es dennoch aber noch Verbesserungsmöglichkeiten gäbe. 

\section{Server}
Unser Administrator Johannes Kirchner gibt bekannt, dass die Inbetriebnahme des neuen Server der Fachschaft ohne Probleme von statten ging. Einige der geplanten VMs laufen bereits Fehlerfrei, darunter zum Beispiel: GitLab, Minecraft (beide noch nicht öffentlich zugänglich) und OwnCloud. 

\section{Sonstiges}
Der Sprecher Georg Schäfer verkündet das er sich beim AStA bezüglich einer Aufwandsentschädigung für die Helfer, welche an der Entwicklung der Satzungsänderung beteiligt waren, erkundigt und diese bewilligt bekommen hat. 

\abstimmung{Es wird über ein geeignetes Kontingent für die Aufwandsentschädigung abgestimmt. Eine Auszahlung von 60 Euro wird als angemessen empfunden.  }{10}{0}{0} %-- {dafür}{dagegen}{enthaltungen}

Außerdem erwähnt der Sprecher, dass die Neuwahlen des Fachschaftsrats, welche im kommenden Semester anstehen, nun nicht mehr auf der Hauptjahresversammlung abgehalten werden müssen, sondern eine eigene Sitzung mit einem dafür gebildeten Wahlausschuss einberufen wird. 
\\\\
\noindent Anschließend tritt unser Kommunikationsminister Marco Lochen von seinem Amt zurück. 
\\\\
\noindent Der Sprecher bedankt sich bei allen Anwesenden für ihr kommen und schließt die Sitzung um 19:45 Uhr. 
\\\\\\
\section{Bestätigung des Protokolls}
Der Sprecher des Fachschaftsrates Informatik sowie der Protokollführer dieses Protokolls bestätigen mit Ihrer Unterschrift unter diesem Protokoll, dass selbiges inhaltlich korrekt ist und alle darin aufgeführten Beschlüsse so wie beschrieben vom Rat getragen werden.
\\

\vspace{3.5cm}
\hrulefill \hfill \hrulefill

\fsiPresident \hfill \protokoller

{\footnotesize (Sprecher)\hfill (Protokollführer)}
\end{document}
